% Options for packages loaded elsewhere
\PassOptionsToPackage{unicode}{hyperref}
\PassOptionsToPackage{hyphens}{url}
\documentclass[
  10pt,
  a4paper,
]{article}
\usepackage{xcolor}
\usepackage[margin=22mm]{geometry}
\usepackage{amsmath,amssymb}
\setcounter{secnumdepth}{-\maxdimen} % remove section numbering
\usepackage{iftex}
\ifPDFTeX
  \usepackage[T1]{fontenc}
  \usepackage[utf8]{inputenc}
  \usepackage{textcomp} % provide euro and other symbols
\else % if luatex or xetex
  \usepackage{unicode-math} % this also loads fontspec
  \defaultfontfeatures{Scale=MatchLowercase}
  \defaultfontfeatures[\rmfamily]{Ligatures=TeX,Scale=1}
\fi
\usepackage{lmodern}
\ifPDFTeX\else
  % xetex/luatex font selection
\fi
% Use upquote if available, for straight quotes in verbatim environments
\IfFileExists{upquote.sty}{\usepackage{upquote}}{}
\IfFileExists{microtype.sty}{% use microtype if available
  \usepackage[]{microtype}
  \UseMicrotypeSet[protrusion]{basicmath} % disable protrusion for tt fonts
}{}
\makeatletter
\@ifundefined{KOMAClassName}{% if non-KOMA class
  \IfFileExists{parskip.sty}{%
    \usepackage{parskip}
  }{% else
    \setlength{\parindent}{0pt}
    \setlength{\parskip}{6pt plus 2pt minus 1pt}}
}{% if KOMA class
  \KOMAoptions{parskip=half}}
\makeatother
\setlength{\emergencystretch}{3em} % prevent overfull lines
\providecommand{\tightlist}{%
  \setlength{\itemsep}{0pt}\setlength{\parskip}{0pt}}
\usepackage{mathrsfs}
\usepackage{mathtools}
\usepackage{xurl}
\emergencystretch=3em
\setcounter{tocdepth}{1}
\newcommand{\Beta}{\mathrm{B}}
\DeclareUnicodeCharacter{00B2}{\ensuremath{^2}}
\DeclareUnicodeCharacter{00B1}{\ensuremath{\pm}}
\DeclareUnicodeCharacter{00D7}{\ensuremath{\times}}
\DeclareUnicodeCharacter{2192}{\ensuremath{\to}}
\DeclareUnicodeCharacter{21D2}{\ensuremath{\Rightarrow}}
\DeclareUnicodeCharacter{21A6}{\ensuremath{\mapsto}}
\DeclareUnicodeCharacter{2208}{\ensuremath{\in}}
\DeclareUnicodeCharacter{2209}{\ensuremath{\notin}}
\DeclareUnicodeCharacter{2212}{\ensuremath{-}}
\DeclareUnicodeCharacter{2227}{\ensuremath{\wedge}}
\DeclareUnicodeCharacter{2228}{\ensuremath{\vee}}
\DeclareUnicodeCharacter{2260}{\ensuremath{\ne}}
\DeclareUnicodeCharacter{2264}{\ensuremath{\le}}
\DeclareUnicodeCharacter{2265}{\ensuremath{\ge}}
\DeclareUnicodeCharacter{2297}{\ensuremath{\otimes}}
\DeclareUnicodeCharacter{00B3}{\ensuremath{^3}}
\DeclareUnicodeCharacter{2074}{\ensuremath{^4}}
\DeclareUnicodeCharacter{2076}{\ensuremath{^6}}
\DeclareUnicodeCharacter{03BB}{\ensuremath{\lambda}}
\DeclareUnicodeCharacter{03BC}{\ensuremath{\mu}}
\DeclareUnicodeCharacter{03B5}{\ensuremath{\epsilon}}
\DeclareUnicodeCharacter{03C4}{\ensuremath{\tau}}
\DeclareUnicodeCharacter{03A6}{\ensuremath{\Phi}}
\DeclareUnicodeCharacter{03B1}{\ensuremath{\alpha}}
\DeclareUnicodeCharacter{03C3}{\ensuremath{\sigma}}
\DeclareUnicodeCharacter{03B4}{\ensuremath{\delta}}
\DeclareUnicodeCharacter{03C0}{\ensuremath{\pi}}
\DeclareUnicodeCharacter{039B}{\ensuremath{\Lambda}}
\DeclareUnicodeCharacter{2202}{\ensuremath{\partial}}
\DeclareUnicodeCharacter{0393}{\ensuremath{\Gamma}}
\DeclareUnicodeCharacter{2205}{\ensuremath{\varnothing}}
\DeclareUnicodeCharacter{221E}{\ensuremath{\infty}}
\DeclareUnicodeCharacter{2200}{\ensuremath{\forall}}
\DeclareUnicodeCharacter{00B9}{\ensuremath{^1}}
\DeclareUnicodeCharacter{00B0}{\ensuremath{{}^\circ}}
\DeclareUnicodeCharacter{211D}{\ensuremath{\mathbb{R}}}
\usepackage{bookmark}
\IfFileExists{xurl.sty}{\usepackage{xurl}}{} % add URL line breaks if available
\urlstyle{same}
\hypersetup{
  hidelinks,
  pdfcreator={LaTeX via pandoc}}

\author{}
\date{}

\begin{document}

{
\setcounter{tocdepth}{3}
\tableofcontents
}
\section{Every time derivative of the actual heat logarithmic
derivative}\label{every-time-derivative-of-the-actual-heat-logarithmic-derivative}

The original family is \(g_t(s)=16H_t(-2i(s-1/2))\), with
\(\partial_tg_t=\tfrac14\partial_s^2g_t\). Its coordinate, analytic
convergence, and complete supported arithmetic receiver are proved in
\href{HEAT_CAUCHY_ARITHMETIC_DERIVATION.md}{HA1--HA26}. This note
computes its time derivatives without replacing any Gamma or endpoint
coefficient by a constant.

\subsection{1. The exact differential-polynomial
recurrence}\label{the-exact-differential-polynomial-recurrence}

Let \(\mathscr P=\mathbb Q[X_0,X_1,\ldots]\) be the polynomial ring in
countably many variables; each polynomial uses only finitely many of
them. Define commuting spatial and evolutionary derivations by \[
D X_j=X_{j+1},\qquad
V X_j=D^j\!\left(\frac{X_2+2X_0X_1}{4}\right).
\tag{HM1}
\] Both extend uniquely by the Leibniz rule and linearity. Their
commutator is a derivation and is zero on every generator, because
\(DVX_j=VDX_j\); therefore \(DV=VD\). Put \[
B_0=X_0,\qquad B_{k+1}=VB_k
=\sum_{j\ge0}\frac{\partial B_k}{\partial X_j}
D^j\!\left(\frac{X_2+2X_0X_1}{4}\right).
\tag{HM2}
\] The sum is finite for each \(k\). Where \(g_t\ne0\), set
\(L_t=g_t'/g_t\). HA3 and the chain rule prove by induction the exact
identity \[
\partial_t^kL_t(s)=B_k(L_t(s),L_t'(s),L_t''(s),\ldots).
\tag{HM3}
\] These are derivatives, not coefficients of a Taylor series; the
Taylor coefficient is \(B_k/k!\). Each formula is valid as a meromorphic
identity in \(s\) and locally holomorphically in \(t\) away from its
zeros.

Give every \(X_j\) polynomial degree one. The derivation \(D\) preserves
this degree. Split \(V=V_1+V_2\), where \[
V_1X_j=\tfrac14X_{j+2},\qquad
V_2X_j=\tfrac12D^j(X_0X_1).
\tag{HM4}
\] Thus \(V_1\) preserves homogeneous degree and \(V_2\) increases it by
one. In particular \(\deg B_k\le k+1\). Its homogeneous component of
degree \(k+1\) is exactly \[
\boxed{B_k^{[k+1]}=\frac1{2^k(k+1)}D^k(X_0^{k+1}).}
\tag{HM5}
\] Here is the full induction. At \(k=0\) the formula gives \(X_0\). As
with \(V\), the derivation \(V_2\) commutes with \(D\). The only term
capable of producing degree \(k+2\) in \(B_{k+1}\) is
\(V_2B_k^{[k+1]}\). Applying the induction hypothesis gives \[
\begin{split}
V_2B_k^{[k+1]}
&=\frac1{2^k(k+1)}D^k\!\left((k+1)X_0^k\frac{X_0X_1}{2}\right)\\
&=\frac1{2^{k+1}(k+2)}D^{k+1}(X_0^{k+2}),
\end{split}\tag{HM6}
\] as required. This polynomial is nonzero; its monomial \(X_0^kX_k\)
has a positive coefficient. Hence \(\deg B_k=k+1\), not only the upper
bound.

For clarity the first two nonconstant polynomials are \[
\begin{split}
B_1&=\frac14(X_2+2X_0X_1),\\
B_2&=\frac1{16}\left(X_4+8X_1X_2+4X_0X_3
+8X_0X_1^2+4X_0^2X_2\right).
\end{split}\tag{HM7}
\] They follow by one and two direct applications of (HM2). In
particular the coefficient of \(X_1X_2\) is eight; no factor is absorbed
into a change of heat time.

\subsection{2. The specified arithmetic
expansion}\label{the-specified-arithmetic-expansion}

On \(\Re s>1\), retain the exact split \[
\begin{split}
L_0(s)&=a(s)+r(s),\\
a(s)&=\frac1s+\frac1{s-1}-\frac12\log\pi+\frac12\psi(s/2),\\
r(s)&=-\sum_{n\ge2}\Lambda(n)n^{-s}.
\end{split}\tag{HM8}
\] The coefficient functions \(a,a',\ldots\) are retained. The expansion
meant here has an explicit rule: in \(B_k\) substitute
\(X_j=a^{(j)}+r^{(j)}\), expand each monomial, and multiply the
absolutely convergent Dirichlet series in the \(r^{(j)}\) factors by
ordered Dirichlet convolution. A term with no \(r\) factor remains its
original coefficient function. This is a specified representation; no
uniqueness claim is made for expansions with arbitrary \(s\)-dependent
coefficients.

Every factor \(r^{(j)}\) has series \[
r^{(j)}(s)=-\sum_{n\ge2}\Lambda(n)(-\log n)^j n^{-s}.
\tag{HM9}
\] Each derivative converges absolutely and locally uniformly on the
half-plane. Finite products do too, because their absolute multiple sums
are bounded by the product of the convergent single sums. Thus the
expansion rule and every spatial derivative used below are justified
there.

Since each nonzero coefficient of one factor is supported on prime
powers, a product of \(j\) such factors is supported on integers with at
most \(j\) distinct prime factors. By (HM4)--(HM5), the arithmetic
coefficient at an integer with more than \(k+1\) distinct prime factors
is zero in this specified expansion of \(\partial_t^kL_t|_0\).

For an integer with exactly \(k+1\) distinct prime factors, write \[
n=\prod_{i=1}^{k+1}p_i^{\alpha_i},\qquad
p_i\ne p_j\ (i\ne j),\quad\alpha_i\ge1.
\tag{HM10}
\] Only the homogeneous component of degree \(k+1\), with every factor
chosen from \(r\), can contribute. Any \(a\) factor would reduce the
number of available prime-power factors. Therefore its coefficient is
the coefficient in \[
\frac1{2^k(k+1)}\partial_s^k(r(s)^{k+1}).
\tag{HM11}
\] To compute it, each of the \(k+1\) ordered factors must receive one
of the distinct prime powers in (HM10). There are exactly \((k+1)!\)
such assignments. No exponent can split between two factors, because
that would leave a different prime without a factor. Hence the
coefficient of \(r^{k+1}\) at \(n\) is
\((-1)^{k+1}(k+1)!\prod_i\log p_i\). Applying \(\partial_s^k\)
multiplies it by \((-\log n)^k\). The two signs combine to minus,
proving \[
\boxed{[n^{-s}]\left.\partial_t^kL_t(s)\right|_0
=-\frac{k!}{2^k}(\log n)^k\prod_{i=1}^{k+1}\log p_i.}
\tag{HM12}
\] The bracket refers exactly to the expansion rule after (HM8). The
displayed coefficient has no remaining \(a(s)\) dependence. At lower
distinct-prime counts all coefficient functions prescribed by that rule
remain present.

For \(k=0\), (HM12) is the original coefficient \(-\log p\) at
\(p^\alpha\). For \(k=1\) it is
\(-\tfrac12\log p\log q\log(p^\alpha q^\beta)\), agreeing with HA19. For
\(k=2\) it is \(-\tfrac12(\log n)^2\log p\log q\log r\) at three
distinct prime factors. Formula (HM12) concerns a time derivative;
dividing by \(k!\) gives the coefficient in the actual Taylor series in
time.

\subsection{3. The global supported
receiver}\label{the-global-supported-receiver}

For the original Cauchy tests at \(z,w\) in \(\Re s>1\), the exact
global family of \href{HEAT_CAUCHY_ARITHMETIC_DERIVATION.md}{HA5--HA20}
gives, for every \(k\), \[
\left.\partial_t^kK_t(z,w)\right|_0
=\frac{B_k(L_0,L_0',\ldots)(w)
+\overline{B_k(L_0,L_0',\ldots)(z)}}{w+\bar z-1}.
\tag{HM13}
\] Indeed the finite-pole formula is holomorphic on a time disc for this
finite test set, so it may be differentiated any finite number of times;
(HM3) supplies the derivatives. All coefficients of \(B_k\) are real.
The fixed supported boundary and every lower coordinate have zero time
derivative for \(k\ge1\). The corresponding geometric-divisor response
is the negative of (HM13) in the top coordinate, with no suppressed
lower distribution, by HA20.

This constructs actual arithmetic mixed-prime data at every
infinitesimal order. The negative coefficient in (HM12) is not the value
of a quadratic form: the matrix denominator, conjugate counterpart,
remaining coefficients, Gamma factors and endpoints in (HM13) are still
present. The complete sign question is the
\href{CAUCHY_WEIL_POSITIVITY_CRITERION.md}{Cauchy--Weil matrix
criterion}.

\subsection{Verification and sources}\label{verification-and-sources}

The accompanying \texttt{check\_\allowbreak{}heat\_\allowbreak{}mixed\_\allowbreak{}prime\_\allowbreak{}recurrence.\allowbreak{}py}
computes the entire recurrence through order six in the differential
polynomial ring, checks its exact degree, and compares every
highest-degree component with (HM5). Its monomial counts are 1, 2, 5,
11, 23, 44 and 82. These finite checks supplement the all-order proof
(HM4)--(HM6).

The human heat source remains Brad Rodgers and Terence Tao,
\href{https://arxiv.org/src/1801.05914v5}{arXiv:1801.05914v5, original
author TeX}, equations \texttt{hoz}, \texttt{htdef} and the subsequent
heat equation. The full supported explicit formula is
\href{https://github.com/KokunoYumeto/zeta-function-research-reader/blob/3c022a0adde0a6aa3d8fc8e43ef42795d88e44b0/workbenches/splitzero-tandem/branches/tau-zero-prime-positivity/SUPPORTED_ZERO_PRIME_WEIL_DERIVATION.md}{SZW19--SZW38},
with its original human sources retained. The logarithmic-derivative
equation is the classical Burgers transformation of the heat equation;
its exact normalization and every recurrence used here have been proved
above. No novelty claim for that transformation is made.

\end{document}
